현대 스마트 제조는 사이버-물리 시스템, 산업용 사물인터넷, 그리고 인공지능을 통합하여 자율적인 생산 환경을 구축하며, 예기치 않은 설비 정지는 심각한 경제적 손실을 초래한다. 이러한 발전에도 불구하고 예지보전 연구는 여전히 단편적이다. 이상 탐지, 잔여수명 예측, 보전 일정 계획은 분리된 문제로 다루어지고, 정적 모델은 센서 드리프트로 인해 성능이 저하되며, 조합적 일정 계획의 복잡성은 기계 수에 따라 빠르게 증가한다. 본 학위논문은 이러한 문제를 단일 실시간 시스템 내에서 해결하기 위한 3계층 통합 아키텍처인 QASAMAP(Quantum-Agentic Smart Manufacturing for Anomaly Detection and Predictive Maintenance) 프레임워크를 제안한다.

제1계층은 50대의 기계에서 69일간 수집된 100{,}000건의 데이터셋에 대한 시공간 예측 및 지도학습을 수행한다. 상호상관 시차 그래프(50개 노드, 538개 유향 엣지)를 구축한 Temporal Graph Convolutional Network(T-GCN)는 5개 센서 채널 전반에 걸쳐 평균 $R^2 = 0.7702$를 달성하여, 독립 LSTM($R^2 = 0.6853$, $+12.4\%$) 및 하이브리드 LSTM--GCN($R^2 = 0.7134$, $+8.0\%$) 베이스라인을 능가한다. 6개 예측 대상에 대해 XGBoost, LightGBM, CatBoost, HistGradientBoosting, TabNet의 포괄적 지도학습 벤치마크를 수행하여 이상 탐지, 고장 유형 분류, 다운타임 위험 및 보전 필요성 예측, 기계 상태, 잔여수명 회귀의 기준 성능을 확립하였다.

제2계층은 Apache Kafka 센서 스트림을 소비하고 기계별로 너비 $W=10$의 윈도우를 버퍼링한 후, 이를 9개의 LLM 기반 에이전트(Monitoring, Correlation, Pattern Learning, Diagnosis, Planning, Action, Self-Reflection, Explainability, Reporting)로 구성된 인지 파이프라인을 통해 단계적 오케스트레이터의 조정 하에 처리하는 4계층 실시간 다중 에이전트 시스템이다. 이질적 LLM 라우팅(GLM-5.1, Kimi-K2.6, DeepSeek-v4-Pro)을 통해 각 에이전트는 작업 특성에 적합한 모델에 매칭된다. 50대 스트리밍 기계에 대한 종단간 검증은 60초 Kafka 주기 내에 전체 파이프라인을 처리한다: 제약 엔진은 50대 중 45대를 CRITICAL 우선순위 수준으로 일정 계획 계층에 전달하고, 후속 MILP 경로는 45대 모든 기계에 대해 일대일 할당 및 용량 제약을 만족하는 일정을 36.9~ms 내에 반환한다.

제3계층은 보전 일정 계획 문제를 혼합 정수 선형 계획법(MILP)과 이차 비제약 이진 최적화(QUBO) 양 형식으로 정식화하고, 동일한 인스턴스에 대해 4가지 솔버 패러다임을 벤치마크한다: 고전적 정확 MILP/CBC, 양자-영감적 진화 알고리즘(QIEA), 모의 양자 어닐링(SQA), 그리고 XY 링 믹서를 사용한 양자 근사 최적화 알고리즘(QAOA). QAOA는 고전적 상태 벡터 시뮬레이션 및 IBM Heron~r2 하드웨어(\texttt{ibm\_fez}, 156 큐비트) 모두에서 실행되었다. 공정한 다중 시드 스케일링 벤치마크(기계 크기당 5개 무작위 시드, 모든 솔버에 동일한 $T$, 동일한 순수 비용 척도)에 따르면, 고전적 MILP/CBC는 $N = 3$부터 $150$까지 모든 시험된 $N$에서 sub-second 해결 시간으로 증명된 최적해를 반환하는 반면, QIEA는 $N=3$에서만 최적해와 일치하고 $N=7$까지는 44--52\% 차이로 실현 가능성을 유지하지만 $N \geq 10$에서는 일대일 할당 실현 가능성을 잃는다. SQA는 페널티 지배적 QUBO 지형으로 인해 모든 시험된 $N$에서 일대일 할당 제약을 충족하지 못한다. QAOA XY 믹서 구성은 일대일 할당 상태를 구조적으로 보존함으로써 시뮬레이션과 하드웨어 실행 모두에서 $N=3$의 MILP 최적해를 복원한다; 하드웨어에서 $N=5$--$10$ 기계의 노이즈 스케일링 경계선은 큐비트 수가 아닌 회로 깊이와 게이트 충실도에 의해 결정된다.

스트리밍 파이프라인에서 직접 도출된 문제 인스턴스에 대해 모든 4개 제약(일대일 기계 할당, $d=4$시간 기간 윈도우에 걸친 슬롯-윈도우 용량, $\rho$-슬롯 위험 시계, 필수 일정 계획 규칙)을 시행한 보완적 복합 제약 실험은 두 가지 뚜렷한 실증적 영역을 드러낸다. $N=3$에서 \texttt{ibm\_fez}의 QAOA-XY($p=2$, 4{,}096 샷, 변환된 깊이 $801$, $644$ CZ 게이트)는 일대일 및 용량 제약 모두를 만족하는 샷의 $10.96\%$ 중 최저 비용 실현 가능 비트스트링을 사후 선택함으로써 MILP 최적해(비용 $= 300.00$)를 정확히 복원한다. $N=5$에서 문제는 MILP/CBC에 대해 구조적으로 실현 불가능하고 실현 가능성 인식 메타휴리스틱(SA, GA, TS, 모두 $4$개 제약 위반에 도달)에 대해서도 실현 불가능하다; QIEA는 비용 $3{,}107.80$(MILP 최선 노력 대비 $+165\%$)에서 느슨하게 실현 가능한 일정을 반환한다. 반대로 QAOA-XY 하드웨어 실행은 비용 $1{,}186.02$(변환된 깊이 $1{,}226$, $1{,}935$ CZ, $0.24\%$ 실현 가능성률)에서 일대일 및 용량 실현 가능한 일정을 사후 선택하며, MILP 최선 노력보다 단지 $1.3\%$ 높고 비교에서 사용 가능한 일정을 반환하는 유일한 솔버이다. 이는 제약 보존 양자 우위이다: XY 링 믹서의 일대일 상태 구조적 보존과 하드웨어 측정 샷에 대한 사후 선택의 결합이, 시험된 모든 고전적 정확, 고전적 휴리스틱, 그리고 QUBO 기반 양자-영감적 방법이 실현 가능성에 실패하거나 $2.6\times$ 비용 프리미엄을 지불하는 영역에서 실현 가능한 일정을 복원한다. 따라서 QASAMAP의 기여는 (i) 60초 실시간 Kafka 예산 내에서 시공간 예측, 다중 에이전트 인지 추론, 그리고 제약 인식 일정 계획에 걸친 정직하고 재현 가능한 벤치마크와 (ii) 제약 보존 양자 최적화가 최선의 고전적 대안에 비해 운영적 가치를 제공하는 영역(MILP 실현 가능성이 무너지는 제약-긴밀 스트리밍 도출 문제)의 구체적 식별이다.
